\documentclass{article}
\usepackage{hyperref}
\usepackage{Style}

\nocite{*} % Comentar si quiero citar
%\addbibresource{bibliografia.bib} % Quitar el comentado si quiero usar bibliografia

\begin{document}

\begin{minipage}{2.5cm}
    \includegraphics[width=2cm]{imagen_puc.jpg}
\end{minipage}
\begin{minipage}{12cm}
    {\sc Pontificia Universidad Católica de Chile\\
    Facultad de Matemáticas\\
    Departamento de Matemática\\
    Profesor: Manuel Cabezas -- Ayudante: Benjamín Mateluna}
\end{minipage}
\vspace{2ex}

{\centerline{\bf Introducción a la Geometría - MAT1304}
\centerline{\bf Ayudantía 9 - Repaso I3}}
\centerline{\bf 9 de abril de 2026}

\vspace{1cm}
\noindent\textbf{Problema 1.} Sea $ABCD$ un paralelogramo y desde $D$ tracemos una recta que 
intersecta al lado $\overline{AB}$ en un punto $E$. Llamemos $F$ al punto de intersección de 
$\overline{DE}$ con $\overline{CB}$. Muestre que
\begin{equation*}
    \frac{\overline{AD}}{\overline{AE}}=\frac{\overline{FB}}{\overline{BE}}
    =\frac{\overline{FC}}{\overline{CD}}
\end{equation*}

\vspace{5mm}
\noindent\textbf{Problema 2.} Sea $\triangle ABC$ escaleno. Sean $P$ y $Q$ en $\overline{AB}$ y
$\overline{AC}$ respectivamente tales que $\overline{PQ}$ es paralela a $\overline{BC}$. La 
circunferencia que pasa por $P$ y es tangente a $\overline{AC}$ en $Q$ intersecta nuevamente a 
$\overline{AB}$ en $R$. Demuestre que $RQCB$ es cíclico.

\vspace{5mm}
\noindent\textbf{Problema 3.} Sea $\triangle ABC$ y $D$ el punto sobre $\overleftrightarrow{AB}$ 
tal que $\overline{CD}$ es la bisectriz externa de $C$, entonces se tiene que
\begin{equation*}
    \frac{\overline{AD}}{\overline{AC}}=\frac{\overline{BD}}{\overline{BC}}
\end{equation*}

\vspace{5mm}
\noindent\textbf{Problema 4.} Sea $O$ el centro de una circunferencia inscrita en $\triangle ABC$.
El rayo $\overrightarrow{AO}$ intersecta a $\overline{BC}$ en $D$. Demuestre que
\begin{equation*}
    \overline{AO}\cdot\overline{BC}=\overline{OD}\cdot(\overline{AB}+\overline{AC})
\end{equation*}

\vspace{1cm}
\noindent\textbf{\underline{Ejercicios Propuestos:}}

\vspace{5mm}
\noindent\textbf{Extra 1.} Dados $\triangle ABC$ y $\triangle DEF$ semejantes con razón $k$, 
pruebe que
\begin{equation*}
    \frac{\abs{\triangle ABC}}{\abs{\triangle DEF}}=k^{2}
\end{equation*}

\vspace{5mm}
\noindent\textbf{Extra 2.} Sea $O$ el centro de una circunferencia circunscrita a un 
$\triangle ABC$ acutángulo. Sea $D$ la intersección de la altura por $A$ con $\overline{BC}$. 
Demostrar que $\angle DAB=\angle OAC$.

%\printbibliography % Quitar el comentado si quiero usar bibliografia

\end{document}